\section{Protocol}
\paragraph{Preregistration and locks.} Sample size, tests, model identities, prompts, code and configuration were frozen before the formal run: the run refused to start unless every locked file hash matched (167 files, six prompt hashes, method versions, model and provider identities, seed commitments and the run environment). The analysis code was locked while campaigns ran and before any formal score existed, and was not changed afterwards; the E3 Reader prompt and the remaining Phase~6 code were locked in the same way before any formal score was read. The test-site master seed and the E3 sample seed were committed with salted hashes in advance. The lock did not include the preregistration documents themselves; their hashes at the lock commit, which was pushed before the run, were recorded afterwards.

\paragraph{Design.} \res{sites} test sites; each language model ran two repeats per site, each a branched Full/Endpoint campaign; the baseline ran two seeds under each condition. The site, not the repeat, is the unit: repeats are averaged within a site. The sample size came from a power analysis on eight pilot sites.

\paragraph{Tests.} The primary family has seven tests with Holm correction: E1, each model's Full final policy minus the baseline's; E2, Full minus Endpoint for each of the four methods. Each is a site-level paired mean difference with a two-sided sign-flip test ($10^5$ flips) and a BCa bootstrap 95\% interval ($10^4$ site resamples); the smallest effect of interest is 2~\eur{}. Secondary results (Benjamini--Hochberg) are each model's gain over its pre-experiment recommendation, each method against the fixed reference, the between-method differences of feedback gains and E3. Format errors, budget exhaustion and fallback recommendations are scored as the method's own result; infrastructure failures could be retried twice and otherwise counted as missing.

\paragraph{Execution.} The run used two rented 128-core servers. A development evaluation repeated on the first server was bit-identical to the pilot desktop's, and evaluations and E3 packets repeated on the second server were bit-identical to the first's. The USD~30 API cap was enforced by provider-side limits summing to it, with pacing reserves in the runner. Deviations: campaigns ran in site order (repeat 0 first) rather than in random order; the simulator is deterministic and all methods of a site ran in the same window. The formal run was completed in two parts. After the first part, an error in the field order of the process predictor's and the baseline's inputs was found; under a new lock, the affected records (all baseline campaigns and \res{repair.branches} language-model Full branches) were rerun from the point of the error, with the baseline radius re-selected on development sites by the preregistered rule. Results are from the corrected data; the original results are in Appendix~\ref{app:original}.

\paragraph{Integrity.} Every planned campaign and evaluation completed, so no data are missing; \res{repair.failed} attempts in the second part that stopped on dropped provider connections were rerun from the same point and are kept; one branch needed more than the preregistered two retries because of a defect in the model client, handled by a recorded decision (Appendix~\ref{app:original}). An audit confirmed one record per planned identity, that all \res{calls.total} model calls (\res{calls.repair} of them in the second part) were served by the planned model and provider with the planned settings, and that the booked spend equals the spend recomputed from per-call records.
